
\SGAProof{Proof of \SGARefLink{sga3:xvii:lemma:3:9:item:i}{3.9(i)}}
The group \(G\) is radicial.  Filtering \(G\) by the \(F^n(G)\), we
reduce to the case in which \(G\) is radicial of height \(1\).
Since \(G\) is nilpotent (\SGARefLink{sga3:xvii:theorem:3-5}{3.5}, \SGARefLink{sga3:xvii:theorem:3:5:item:i}{(i)}\(\Rightarrow\)\SGARefLink{sga3:xvii:theorem:3:5:item:iii}{(iii)}), and since
the center of an algebraic group is representable
(\SGARefLink{sga3:viii:examples:6:5:item:e:01}{Exposé VIII, 6.5(e)}), we may consider the ascending central series
of \(G\).  This series is plainly characteristic in \(G\), and
reduces us to the case in which \(G\) is also commutative.

Put \(\mathfrak g=\operatorname{Lie}(G)\).  The \(p\)-th-power
morphism \(\pi\) is then additive on \(\mathfrak g\)
(Exposé VII\textsubscript{A}); we shall reduce to the case in which
it is zero.  For every prescheme \(S\) over \(k\), set
\[
  \mathfrak g_S=\mathfrak g\otimes_k\mathcal O_S,
\]
and let \(\mathfrak h_S\) be the subsheaf of abelian groups of
\(\mathfrak g_S\) that is the image of \(\mathfrak g_S\) under
\(\pi_S\).  Finally, let \(\overline{\mathfrak h}_S\) be the
subsheaf of \(\mathcal O_S\)-modules of \(\mathfrak g_S\) generated
by \(\mathfrak h_S\).  It is clear that
\[
  \overline{\mathfrak h}_S
  =\overline{\mathfrak h}_k\otimes_k\mathcal O_S
\]
and that \(\overline{\mathfrak h}_S\)\footnote[4]{Editor's note:
\(\mathfrak h\) has been changed to \(\overline{\mathfrak h}\).}
is a characteristic \(p\)-Lie subalgebra of \(\mathfrak g_S\);
that is, it is stable under the \(S\)-functor
\(\operatorname{Aut}_{p\text{-}\mathrm{Lie}}(\mathfrak g)\).
It follows from Appendix~II, \SGARefLink{sga3:ii:subsection:2-2}{2.2} that
\(\overline{\mathfrak h}_k\) is the Lie algebra of an algebraic
subgroup \(H\) of \(G\) that is characteristic in \(G\).

Moreover, in view of \SGARefLink{sga3:xvii:corollary:3:7}{3.7} and \SGARefLink{sga3:xvii:lemma:3:9:bis}{Lemma 3.9 bis} below, if
\(G\ne\{e\}\), then \(H\ne G\), since
\(\overline{\mathfrak h}_k\ne\mathfrak g\).  On the other hand, if
\(G/H=G''\), then
\[
  \operatorname{Lie}G''
  =\mathfrak g/\overline{\mathfrak h}_k
  \qquad\text{(Appendix~II, 2.2),}
\]
and consequently the \(p\)-th power is zero in
\(\operatorname{Lie}G''\).  Induction on
\(\dim\operatorname{Lie}G\) therefore reduces us to the case in
which \(\operatorname{Lie}G\) is a \(p\)-Lie algebra whose
\(p\)-th power is zero.  But then \(\operatorname{Lie}G''\) is
isomorphic to \(\operatorname{Lie}((\boldsymbol\alpha_p)^r)\) for
some suitable integer \(r\geq0\) (Appendix~II, \SGARefLink{sga3:xvii:proposition:2:1}{2.1}), and hence
\(G''\) is isomorphic to \((\boldsymbol\alpha_p)^r\)
(Appendix~II, \SGARefLink{sga3:ii:subsection:2-2}{2.2}).  It remains to prove the following lemma.

\SGAStatement{Lemma 3.9 bis}{Let \(k\) be a field of characteristic
\(p>0\), let \(\mathfrak g\) be a commutative unipotent
\(p\)-Lie algebra (\SGARefLink{sga3:xvii:definition:3:6}{3.6}) of finite dimension over \(k\), and let
\(\mathfrak h\) be the \(p\)-Lie subalgebra of \(\mathfrak g\)
generated by the image of the \(p\)-th power in \(\mathfrak g\).
If \(\mathfrak g\ne0\), then \(\mathfrak h\ne\mathfrak g\).}
\SGARefTarget{sga3:xvii:lemma:3:9:bis}

Indeed, since \(\mathfrak g\) is commutative, \(\mathfrak h\) is
simply the \(k\)-vector subspace of \(\mathfrak g\) generated by the
\(X^{(p)}\), \(X\in\mathfrak g\).  If \(\mathfrak g\ne0\) and is
unipotent, there exists an \(X\in\mathfrak g\), \(X\ne0\), such
that \(X^{(p)}=0\).  Let \(X_1,\ldots,X_n\) be a basis of a
complement in \(\mathfrak g\) to the line \(kX\).  The Lie algebra
\(\mathfrak h\) is then the \(k\)-vector subspace of
\(\mathfrak g\) generated by
\[
  X_1^{(p)},\ldots,X_n^{(p)};
\]
it therefore has dimension at most
\(n=\dim\mathfrak g-1\).

\SGAProof{Proof of \SGARefLink{sga3:xvii:lemma:3:9:item:ii}{3.9(ii)}}
The group \(G\) is smooth and connected.  In this case the
descending central series of \(G\) is represented by smooth,
connected, characteristic algebraic subgroups \(G_i\)
(Exposé VI\textsubscript{B}, \SGARefLink{prop:sga3-VIB-8.3}{8.3} and \SGARefLink{prop:sga3-VIB-7.4}{7.4}), and \(G_i=0\) for
sufficiently large \(i\), since \(G\) is nilpotent
(\SGARefLink{sga3:xvii:theorem:3-5}{3.5}, \SGARefLink{sga3:xvii:theorem:3:5:item:i}{(i)}\(\Rightarrow\)\SGARefLink{sga3:xvii:theorem:3:5:item:iii}{(iii)}).  It is enough to prove \SGARefLink{sga3:xvii:lemma:3-9}{3.9} for the
groups \(G_i/G_{i+1}\), which reduces us to the case in which \(G\)
is also commutative.

For every integer \(n>0\), let \(G_n\) be the algebraic subgroup of
\(G\) that is the image of \(G\) under the \(p^n\)-th-power
morphism.  The group \(G_n\) is therefore smooth, connected, and
characteristic; and it follows from \SGARefLink{sga3:xvii:definition:1:1}{Definition 1.1} of unipotent
groups that \(G_n=0\) for sufficiently large \(n\).  Replacing \(G\)
by \(G_n/G_{n+1}\), we may further suppose that \(G\) is annihilated
by raising to the \(p\)-th power.  But then, by J.-P. Serre,
\emph{Groupes algébriques et corps de classes}, Chap.~VII,
Prop.~11, \(G\) is a form of \((\mathbb G_a)^r\) for a suitable
integer \(r\).
